% ============================================================
% f2014_guangdong_science_q18_pyramid.tex
% 2014 广东卷（理科）第 18 题：四边形 ABCD 为正方形，
% PD⊥平面 ABCD，∠DPC=30°，AF⊥PC 于 F，FE∥CD 交 PD 于 E。
% 来源 PDF：普通高考/2014/2014广东理.pdf
% 原页：第 1 页（A3 横排）第 3 栏下部，印面 877，第 18 题题干下方
% 说明：底面 ABCD 画成平行四边形以示意正方形；PD 为自 D 向下的
% 棱，E 在 PD 上，F 在 PC 上，EF∥CD；被底面遮挡的 AD、DC、DF
% 用虚线，其余可见棱用实线。
% ============================================================
\begin{tikzpicture}[>=Stealth, line width=0.7pt, font=\small,
  x=1cm, y=1cm, line cap=round, line join=round]

  \coordinate (A) at (0,3.4);
  \coordinate (B) at (4.2,3.4);
  \coordinate (C) at (5.2,0.9);
  \coordinate (D) at (1.0,0.9);
  \coordinate (P) at (-0.4,-1.8);
  \coordinate (E) at (0.65,0.225);
  \coordinate (F) at (3.8,0.225);

  % 被遮挡的棱
  \draw[dashed] (A) -- (D);
  \draw[dashed] (D) -- (C);
  \draw[dashed] (D) -- (F);

  % 可见的棱
  \draw (A) -- (B) -- (C);
  \draw (P) -- (E) -- (D);
  \draw (P) -- (C);
  \draw (E) -- (F);
  \draw (A) -- (C);
  \draw (A) -- (F);

  \fill (A) circle (1.1pt) node[above left] {$A$};
  \fill (B) circle (1.1pt) node[above right] {$B$};
  \fill (C) circle (1.1pt) node[right] {$C$};
  \fill (D) circle (1.1pt) node[above right, inner sep=2pt] {$D$};
  \fill (P) circle (1.1pt) node[below left] {$P$};
  \fill (E) circle (1.1pt) node[left] {$E$};
  \fill (F) circle (1.1pt) node[below right] {$F$};
\end{tikzpicture}
